\section{The ruler and cell transfer}\label{sec:ruler}
We state precisely the imported geometric interface. With
$r=1/16$, a perfect cell of pitch $p\geq2$ has interior side $h=p-1$,
and the classical density is
\begin{equation}\label{eq:ruler}
 f_g(p)=-\frac gp+\frac r{p^2},\qquad
 g_p=r\left(\frac1p+\frac1{p+1}\right),\qquad
 f_g(p)-f_g(p+1)=\frac{g_p-g}{p(p+1)}.
\end{equation}
The five symbols are $W,V,A,S,D$. For orientation, the allowed horizontal
pairs are
\[
 WW,AA,SA,DS,DD,AV,SV,VS,VD;
\]
the allowed top-to-bottom pairs are
\[
 WS,WA,WV,SW,DW,VW,AA,AS,SD,DD,VV.
\]
Every other pair costs one. An onsite $W$ bonus is $-g$; a vertical
$VW$ bond costs $r$. We use the forcing and charging theorems of
\cite[Lemmas 2.1--2.2]{gap}, rather than a new rule-level proof.

For small $g>0$, $p_*(g)=p$ on $(g_p,g_{p-1})$, with both $p,p+1$
minimizing at $g_p$. Each stripe can choose its own horizontal phase.
The source's Corollary 2.3 gives separation from every nonminimizing
classical block by at least $\min(5/8,L\delta(g))$, after also making
the vacuum excess exceed that bound, where
\[
 \delta(g)=\min_{p\notin\mathcal P_*(g)}
       p(f_g(p)-\min_q f_g(q))>0.
\]
At a tie $\mathcal P_*(g)$ contains both minimizers. For fixed rational
$g$, a factorial-volume cutoff is computable. Every fixed positive
$g$ also has eventual commensurability and separation. In particular
\begin{equation}\label{eq:pitch}
 p_*(g)=2r/g+O(1),\qquad \log_2h_*(g)=\log_2(1/g)+O(1).
\end{equation}
For $g\leq0$ the two wall-free vacua minimize, with classical margin
at least $r$.

The flags in \cite[Sections 5 and 11]{gap} tie the top, bottom, left
and right of a tableau to its enclosing walls. In a correct cell the
classical content is the unique functional tableau of an $h\times h$
box. All extra tests are nonnegative unit penalties. All quantum terms
below are positive, conditioned on local classical labels, and act
either onsite or on a bond entirely inside a cell. Exterior quantum
registers have one-dimensional frozen kernels and a specified positive
excitation bound.

\begin{lemma}[Reparametrization]\label{lem:unit}
Suppose $\phi(g)$ samples every sufficiently large integer, only large
integers as $g\downarrow0$, and $\ln\phi(g)=\ln(1/g)+O(1)$ uniformly.
If $Q(h)>0$ and $|\ln Q(h)|=O(\ln h)$, the lower and upper power
exponents of $Q(\phi(g))$ are those of $Q(h)$. If $Q(h)=O(\ln h)$,
the same is true of the logarithmic coefficients.
\end{lemma}
\begin{proof}
On a plateau with $\phi(g)=h$, the denominators differ by $O(1)$.
The bounded normalized numerators make the difference of ratios $o(1)$.
Every tail integer occurs, and small $g$ excludes any finite prefix.
Consequently the two sets of tail limit points coincide. The argument
does not require a computable convergence rate.
\end{proof}

\begin{proposition}[Selected-block transfer]\label{prop:transfer}
Suppose a selection theorem gives minimizing perfect blocks of pitch
$p$, with all other classical blocks at energy at least $E_0+\gamma$.
Let the frozen excitation bound be $\gamma_f>0$, and include all
tableau, flag, and exterior margins in
$\gamma_{\rm eff}=\min(\gamma,\gamma_f,\gamma_{\rm other})$.
In each minimizing block assume exactly
\begin{equation}\label{eq:tensorsum}
 H=E_0I+\sum_c(K_h-e(h)I)_c+\sum_w F_w,
 \qquad e(h)=\min\spec K_h,
\end{equation}
on disjoint factors, with $\ker F_w=\C\ket{0}_w$.
Below $E_0+\gamma_{\rm eff}$, the full spectrum is the sum spectrum
of the cells, with multiplicities summed over the minimizing blocks.
The full gap equals $\Delta^+(K_h)$ if that number is less than
$\gamma_{\rm eff}$; otherwise the full gap is at least
$\gamma_{\rm eff}$. Here $\Delta^+(K_h)=+\infty$ if the cell has
no positive excitation above its ground space.
The ground space is
\begin{equation}\label{eq:groundspace}
 \bigoplus_{\sigma\ {
m minimizing}}\ket{\sigma}\otimes
       \left(\bigotimes_c\ker(K_h-e(h)I)\right)\otimes\ket{\rm frozen}.
\end{equation}
Effective specified cell quantities with finite normalized endpoints
transfer by Lemma~\ref{lem:unit}. Uniformly computable scalar cell data
have lower, upper, and existing finite limits in
$\LS,\LP,\Dtwo$, respectively.
\end{proposition}
\begin{proof}
All classical labels are conserved. Nonminimizing blocks and excited
frozen factors cannot contribute below the stated margin. Diagonalize
the disjoint factors in \eqref{eq:tensorsum}; eigenvalues add and
ground spaces tensor. This proves the spectral statement and
\eqref{eq:groundspace}. If $d_h=\dim\ker(K_h-e(h)I)$, the degeneracy
is $N_{\min}d_h^{L^2/p^2}$. The stripe phase count has
$\log N_{\min}=O((L/p+1)\log p)$ at fixed pitch, so its entropy density
vanishes; the quantum residual entropy density is $(\log d_h)/p^2$.

Reparametrization proves the transfer assertion for the specified
quantity. For arithmetic, uniformly approximate each normalized cell
datum by a rational within $2^{-h}$. For a rational sequence $a_h$,
\[
 q<\liminf_h a_h\quad\Longleftrightarrow\quad
 \exists r\in\Q\ \exists H\ [r>q\ \text{and}\ \forall h\geq H:\ a_h\geq r].
\]
This is a second-level existential left cut. The upper endpoint is
dual; an existing finite limit is computable in the halting set.
For an irrational endpoint complement the appropriate strict cut;
a rational endpoint is itself computable. This is membership of an
individual real in a cut class, not a uniform exact-comparison
algorithm for arbitrary Cauchy names. It agrees with
\cite[Proposition 10.1]{gap} and the prior calculus \cite{zheng}.
\end{proof}

If every cell ground state is unique, a single-block ground vector is
a product over cells. More generally, that product is an available
choice. For such states, restrictions to any region tensor over its
cell intersections, so all R\'enyi entropies add over cut cells;
whole-cell unions have zero entropy and different cells' observables
factor. These conclusions are not statements about every vector in
\eqref{eq:groundspace}. A free logical qubit in each cell permits a
Bell state across two cells in a single classical block, giving a
whole cell entropy one. Repeated finite Bell dimers even give a pure
clustering example. Arbitrary cell marginals must therefore be used
when cell kernels are degenerate.

\begin{lemma}[Small-energy perfect-stripe staircase]\label{lem:shifted}
Let $0\leq e(h)\leq ch^{-3}$ for $h\geq1$ and set
\[
 F_g(p)=-g/p+(r+e(p-1))/p^2,\qquad p\geq2.
\]
Define $d_p=((p+1)/p)e(p-1)-(p/(p+1))e(p)$ and $g'_p=g_p+d_p$.
Then $F_g(p)-F_g(p+1)=(g'_p-g)/(p(p+1))$.
The ties $g'_p$ strictly decrease for
$p\geq p_0=\lceil4c/r\rceil+2$, and all sufficiently large pitches
are selected in order among perfect stripes, with $p_*g\to2r$.
\end{lemma}
\begin{proof}
Subtract the two densities. Nonnegativity gives
$|d_p|\leq2c/(p-1)^3$, so
\[
 g'_p-g'_{p+1}\geq\frac{2r}{p(p+2)}-\frac{4c}{(p-1)^3}.
\]
For $p\geq p_0$ use $4c\leq r(p-2)$ and
$2(p-1)^3-p(p+2)(p-2)=p^3-6p^2+10p-2>0$ for integers $p\geq2$.
The last polynomial is positive at $2,3,4$ and increasing thereafter;
the case $c=0$ is immediate. At sufficiently small $g$, the finitely
many smaller pitches have positive density, whereas a pitch near
$2r/g$ has negative density. Adjacent differences then change sign
exactly once on the tail. Finally $g'_p=2r/p+O(p^{-2})$.
\end{proof}

This lemma compares perfect stripes. Selection against irregular
configurations with positive cell energy requires the local charging
hypothesis in Section~\ref{sec:conditional}. Zero cell energy needs
only the imported ruler theorem.
